%% kapa: exam=1 round=34 year=2023 period=1 form=A subject=civil_law subject_ko=민법 seq=30 number=30
%% source: Q-Net 감정평가사 기출문제 (cst00309) artlSeq=5211789 file=raw/1008/5211789/2023년도 제34회 감정평가사 1차시험 1교시.pdf
%% conversion: kordoc 4.24.51 -> tools/gen1.py (auto); review=proofread-sonnet; flags=none
\begin{question}{30}
부합에 관한 설명으로 옳지 않은 것은? (다툼이 있으면 판례에 따름)
\begin{choices}
  \item 부동산에의 부합 이외에 동산 상호간의 부합도 인정된다.
  \item 동산 이외에 부동산은 부합물이 될 수 없다.
  \item 동일인 소유의 부동산과 동산 상호간에는 원칙적으로 부합이 인정되지 않는다.
  \item 분리가 가능하지만 분리할 경우 상호 부착되거나 결합된 물건의 경제적 가치가 심하게 손상되는 경우에도 부합이 인정된다.
  \item 부동산의 소유자는 원칙적으로 그 부동산에 부합한 물건의 소유권을 취득한다.
\end{choices}
\end{question}
\answer{2}
